% cover_letter.tex
% IEEE Transactions on Power Delivery — submission cover letter
% Paper: "IBR-Adaptive Line Differential Protection: Harmonic Discrimination,
%         Sequence Elements, and Transient Differential Security"
% Authors: A. V. Eluvathingal and K. S. Swarup, IIT Madras
\documentclass[11pt,a4paper]{article}
\usepackage[T1]{fontenc}
\usepackage[utf8]{inputenc}
\usepackage[margin=2.5cm]{geometry}
\usepackage{microtype}
\usepackage{parskip}
\usepackage{newtxtext}
\pagestyle{empty}

\begin{document}

\noindent
\today

\bigskip

\noindent
Editor-in-Chief\\
\textit{IEEE Transactions on Power Delivery}

\bigskip

\noindent
Dear Editor-in-Chief,

\medskip

We submit for your consideration the original research article:

\begin{center}
\textbf{``IBR-Adaptive Line Differential Protection:\\
Harmonic Discrimination, Sequence Elements, and Transient Differential
Security''}\\[4pt]
\textit{Anoop V. Eluvathingal and K. Shanti Swarup}\\
Department of Electrical Engineering, Indian Institute of Technology Madras,
Chennai~600~036, India
\end{center}

\medskip\noindent\textbf{Motivation and problem addressed.}
The progressive displacement of synchronous generators by inverter-based
resources (IBRs) --- solar PV, battery storage, and grid-forming wind
converters --- imposes a well-documented but unsolved problem on line
differential (87L) protection.
IBR fault-current limiters restrict the short-circuit contribution to
1.0--1.2\,pu, producing a differential signal of 0.06--0.20\,pu at the
relay CT primary, which falls below the standard 87L pickup threshold.
In our deterministic test matrix, the baseline relay fails on 9 of 21
scenarios (57\% failure rate) at low IBR penetration.
Three independent failure modes are responsible: magnitude suppression for
three-phase faults, harmonic-induced confidence collapse, and
negative-sequence suppression for asymmetric faults.
No published 87L scheme addresses all three simultaneously, and no prior
work has exploited the fault-inception transient as a primary detection
feature for symmetric IBR faults.

\medskip\noindent\textbf{Contributions.}
This paper proposes a five-element IBR-adaptive 87L protection chain that
closes all three gaps:
\begin{enumerate}
  \item \textbf{Physics-constrained adaptive threshold (C1):}
        The bias slope $k_m^\ast$ is derived analytically from simultaneous
        CT-mismatch stability and high-impedance-fault sensitivity constraints
        as a closed-form function of the real-time IBR current ratio $k_{\rm ibr}$.
        This eliminates the conservative fixed setting that previously caused
        either missed trips (low $k_{\rm ibr}$) or false trips (high
        $k_{\rm ibr}$).

  \item \textbf{Frequency-adaptive harmonic pre-filter (C2):}
        A zero-crossing frequency tracker maintains the pi-section charging
        canceller orthogonal to 47--53\,Hz; a QR-based harmonic regressor
        removes IBR-injected odd harmonics, restoring the relay confidence
        score from $c = 0.37$ to $c \geq 0.99$ under 10\,\% THD.

  \item \textbf{Negative- and zero-sequence differential elements (C3):}
        An 87LN element provably achieves differential-to-bias ratio
        $\rho_2 = 2$ irrespective of $k_{\rm ibr}$ when the IBR suppresses its
        negative-sequence output --- inverting the conventional assumption that
        IBR sequence suppression is always detrimental.
        An 87LN0 element extends coverage to solidly-earthed pure-IBR earth
        faults.

  \item \textbf{Transient differential current signature element (C4):}
        TDCS-87L exploits the step-change in differential RMS at fault
        inception, which does not occur for external faults or load steps.
        An analytical bound proves detection for all faults with
        $k_{\rm ibr} \geq 0.058$\,pu.
        Monte Carlo Group~C (5\,000 external fault trials) confirms a
        15.5\,\% security margin against Class~5P CT mismatch at the
        joint worst-case parameter corner.

  \item \textbf{Solar PV two-parameter model and DDR selector (C5):}
        The zero-DC-offset property of GFL PV fault current is exploited to
        reduce the zone estimator from four parameters (condition number
        $\kappa_n \in [3,36]$) to two ($\kappa_n = 1$ analytically).
        A DC Decay Ratio gate routes each relay cycle to the correct
        estimator, reducing the median PV-fault trip time from 63\,ms to
        42\,ms.
\end{enumerate}

\medskip\noindent\textbf{Validation.}
The chain is validated deterministically on 30 scenarios (21 SG/IBR + 9 PV),
with zero regression on the 10 synchronous-generator reference cases, and
probabilistically on 20\,000 Monte Carlo trials ($N = 5\,000$ per group,
four groups).
Groups~A--C (three-phase internal, SLG internal, external faults) achieve
TPR\,=\,1.000 and FPR\,=\,0.000 with 95\,\% Wilson CI lower bound
$\geq 0.9992$.
Group~D (PV internal faults) achieves TPR\,=\,0.9984
(CI$_{\rm lo} \geq 0.9976$).
All results exceed IEC~60255-181 performance targets.

\medskip\noindent\textbf{Fit for \textit{IEEE Transactions on Power Delivery}.}
The paper addresses a fundamental protection challenge that is directly
relevant to the Transactions' readership: relay engineers and system
planners dealing with the ongoing transition to high-IBR networks.
The combination of closed-form analytical derivation, formal propositions
with proofs, and large-scale Monte Carlo certification meets the technical
depth expected of a Transactions paper.
The work is original and has not been submitted elsewhere.

\medskip\noindent\textbf{Suggested reviewers.}
Given the subject matter, we suggest the following independent experts (none
have been involved in this work):
\begin{itemize}
  \item Prof.\ Ali Hooshyar, University of Toronto (IBR protection,
        line differential relaying)
  \item Prof.\ Evangelos Farantatos, EPRI (IBR fault current characterisation,
        adaptive protection)
  \item Prof.\ Hossein Hooshyar, NTNU (IBR-impacted protection coordination)
  \item Dr.\ Bogdan Kasztenny, SEL (differential relay design and
        negative-sequence elements)
\end{itemize}

\medskip\noindent\textbf{Authorship and conflicts.}
Both authors declare no conflicts of interest.
The contributions statement is provided separately in the portal upload.
All data and simulation parameters are disclosed in the paper body and
referenced technical reports; simulation code is available upon request.

\medskip

We look forward to the editorial review and are available to respond to
any reviewer queries.

\bigskip

\noindent Yours sincerely,

\bigskip\bigskip

\noindent
\textbf{Anoop V. Eluvathingal} \hfill \textbf{K. Shanti Swarup}\\
PhD Scholar (Corresponding Author) \hfill Professor\\
Department of Electrical Engineering \hfill Department of Electrical
Engineering\\
IIT Madras, Chennai 600 036, India \hfill IIT Madras, Chennai 600 036,
India\\
\texttt{ee21d017@smail.iitm.ac.in} \hfill
\texttt{swarup@ee.iitm.ac.in}

\end{document}
